% Part: first-order-logic
% Chapter: axiomatic-deduction
% Section: deduction-theorem

% verification of properties of provability needed for maximally
% consistent sets in the completeness chapter.

\documentclass[../../../include/open-logic-section]{subfiles}

\begin{document}

\iftag{FOL}
      {\olfileid{fol}{axd}{ded}}
      {\olfileid{pl}{axd}{ded}}

\olsection{مبرهنة الاستنباط}

قد ظهر أن كتابة ال!!{derivation}s في النسق البديهي ذات مؤونة،
وأن العثور على ال!!{derivation}s قد يعسر. وغالبًا ما يكفينا،
بدل سرد متتاليات طويلة من ال!!{formula}s، أن نثبت وجود
ال!!{derivation}s بنتائج يسيرة. وقد تقدّم منها ثلاث:
فبـ\olref[ptn]{prop:reflexivity} نحكم بـ~$\Gamma
\Proves !A$ متى كان~$!A \in
\Gamma$؛ وبـ\olref[ptn]{prop:monotonicity} يلزم من~$\Gamma \Proves !A$
أن~$\Gamma \cup \{!B\} \Proves !A$؛
وبـ\olref[ptn]{prop:transitivity} يلزم من~$\Gamma \Proves !A$
و$!A \Proves !B$ أن~$\Gamma \Proves !B$.
ونضم إليها الآن صورةً من إثبات المقدَّم تتعلق بوجود الاشتقاقات
نفسه:


\begin{prop}
  \ollabel{prop:mp} إذا كان $\Gamma \Proves !A$ و$\Gamma \Proves !A \lif
  !B$، فإن $\Gamma \Proves !B$.
\end{prop}

\begin{proof}
  يكفي أولًا أن نلاحظ~$\{!A, !A \lif !B\} \Proves !B$:
  \begin{derivation}
    1. & $!A$ & \Hyp\\
    2. & $!A \lif !B$ & \Hyp\\
    3. & $!B$ & 1، 2، إثبات المقدَّم
  \end{derivation}
  وبحسب \olref[ptn]{prop:transitivity}، نحصل على $\Gamma \Proves !B$.
\end{proof}

وأعظم ما ننتفع به في هذا الباب مبرهنة الاستنباط:

\begin{thm}[مبرهنة الاستنباط]
\ollabel{thm:deduction-thm}
\iftag{FOL}{متى كان~$\Gamma \cup \{!A\} \Proves !B$ باشتقاق لا يعمل
إلا بالفروض والبديهيات وقاعدة إثبات المقدَّم، كان
$\Gamma \Proves !A \lif !B$. وفي الجهة الأخرى، متى كان
$\Gamma \Proves !A \lif !B$، كان~$\Gamma \cup \{!A\} \Proves !B$
من غير هذا القيد.}{تكون~$\Gamma \cup \{!A\} \Proves !B$ إذا وفقط
إذا كان~$\Gamma \Proves !A \lif !B$.}
\end{thm}

\begin{proof}
أما اتجاه «إذا»، فمن~$\Gamma \Proves !A \lif !B$ يلزم
$\Gamma \cup \{!A\} \Proves !A \lif !B$ بمقتضى
\olref[ptn]{prop:monotonicity}. ولدينا~$\Gamma \cup \{!A\} \Proves !A$
بمقتضى \olref[ptn]{prop:reflexivity}. فبـ\olref{prop:mp} نحصل على
$\Gamma \cup \{!A\} \Proves !B$.


وأما اتجاه «فقط إذا»، فنستقرئ على طول !!{derivation} الصيغة
$!B$ من~$\Gamma \cup \{!A\}$.

في الأساس نأخذ !!{derivation} طوله~$1$.
فال!!^a{derivation} لـ~$!B$ من~$\Gamma \cup \{!A\}$ إذا كان
طوله~$1$ لم يشتمل إلا على~$!B$. ولا تصح هذه الخطوة إلا إذا
كان~$!B \in \Gamma \cup \{!A\}$ أو كانت بديهية.
فإن كان~$!B \in \Gamma$ أو كانت بديهية، حصل~$\Gamma \Proves !B$.
وبالبديهية~$\Gamma
\Proves !B \lif (!A \lif !B)$، وهي \olref[prp]{ax:lif1}،
تعطينا \olref{prop:mp} النتيجة~$\Gamma \Proves !A \lif !B$.
وإن كان~$!B \in \{ !A\}$، صح~$\Gamma \Proves !A \lif !B$؛
إذ ال!!{formula}~$!A \lif !B$ هي~$!A \lif !A$،
وقد أمكن !!{derive} هذه في \olref[pro]{ex:identity}.

وفي الخطوة نأخذ !!a{derivation} لـ~$!B$ من
$\Gamma \cup \{!A\}$ خاتمته~$!B$ مبرّرة بإثبات المقدَّم.
فإن لم تكن كذلك، كان~$!B \in \Gamma$ أو~$!B
\ident !A$ أو كانت~$!B$ بديهية؛ فيجري عليها تعليل الأساس.
أما إذا كانت بإثبات المقدَّم، فقد سبقت في ال!!{derivation}
$!C \lif !B$ و$!C$ لبعض ال!!{formula}~$!C$؛ أي
$\Gamma \cup \{!A\} \Proves !C \lif !B$ و$\Gamma \cup \{!A\}
\Proves !C$. وال!!{derivation}s المنتهية بهاتين الخطوتين أقصر،
فيشملها فرض الاستقراء، ويحصل منه:
\begin{align*}
  & \Gamma \Proves !A \lif (!C \lif !B); \\
  & \Gamma \Proves !A \lif !C.
\end{align*}
ومن جهة أخرى، لدينا
\[
\Gamma \Proves (!A \lif (!C \lif !B)) \lif
((!A\lif !C)  \lif (!A \lif !B)),
\]
بمقتضى \olref[prp]{ax:lif2}. نطبق \olref{prop:mp} مرتين،
فنصل إلى~$\Gamma \Proves !A \lif !B$، وهو المطلوب.
\end{proof}

ومن هذا يظهر وجه اختيار \olref[prp]{ax:lif1}
و\olref[prp]{ax:lif2}؛ فقد صيغتا بحيث تصح مبرهنة الاستنباط.

ونختم بنتائج نافعة في !!{derivability}، ونجعل إثباتها تمارين.

\begin{prop}
  \ollabel{prop:derivfacts}
\begin{enumerate}
\item $\Proves (!A \lif !B) \lif ((!B \lif !C)
  \lif (!A \lif !C))$؛ \ollabel{derivfacts:a}
\item إذا كان $\Gamma \cup \{ \lnot !A\}
  \Proves \lnot !B$، فإن $\Gamma \cup \{ !B\} \Proves
  !A$ (المقابلة)؛ \ollabel{derivfacts:b}
\item  $\{ !A, \lnot!A\} \Proves
    !B$ (من الباطل يلزم أي شيء، الانفجار)؛ \ollabel{derivfacts:c}
\item  $\{ \lnot\lnot!A\} \Proves
  !A$ (حذف النفي المزدوج)؛\ollabel{derivfacts:d}
\item إذا كان $\Gamma \Proves \lnot\lnot!A$، فإن $\Gamma \Proves
  !A$؛\ollabel{derivfacts:e}
\end{enumerate}
\end{prop}

\tagprob{FOL}
\begin{prob}
  أثبت \olref[fol][axd][ded]{prop:derivfacts}.
\end{prob}
\tagendprob

\tagprob{notFOL}
\begin{prob}
  أثبت \olref[pl][axd][ded]{prop:derivfacts}.
\end{prob}
\tagendprob

\end{document}
